% Л15 -- Тестирование и качество. Slides.
%
% Companion to the spoken script in ../L15-testing.md, which is the source of
% truth: the slides carry diagrams, code, real terminal output and short
% marks, and the delivery lives in the script only -- this deck carries no
% speaker notes. Slide text is in Russian because that is the content;
% everything else in this file is English.
%
% Build:  make          (latexmk -pdf; engine: pdflatex)
%
% The C# listings were compiled on dotnet SDK 10.0 in scratch projects
% (`HelloWorld` + `HelloWorld.Tests`, a `Benchmarks` console project and the
% lab reference solution); terminal output on the slides is captured from
% runs of those projects. See ../CLAUDE.md for the verification table.
%
% House rules: ../CLAUDE.md and ../../../CLAUDE.md. Theme: ../../../../theme/.

\documentclass[aspectratio=169,10pt,t]{beamer}

\usetheme{dotnet}
\usetikzlibrary{arrows.meta,positioning,calc,decorations.pathreplacing,fit}
\usepackage{array}   % ragged-right table columns

% A ragged-right paragraph column: justified cells in narrow columns hyphenate
% badly, and a split word is the one thing the back row cannot read.
\newcolumntype{L}[1]{>{\raggedright\arraybackslash}p{#1}}

\dotnetlecture{Л15}

\title[Л15. Тестирование и качество]{Тестирование и качество}
\author{Дмитрий Курочкин}
\date{}

% Code sizes, one step down from \footnotesize. Smaller than 6.6pt does not
% carry to the back of the room.
\newcommand{\codesm}{\ttfamily\fontsize{7.6}{9.0}\selectfont}
\newcommand{\codexs}{\ttfamily\fontsize{7.0}{8.3}\selectfont}
\newcommand{\codexxs}{\ttfamily\fontsize{6.8}{7.8}\selectfont}
\newcommand{\termsm}{\ttfamily\fontsize{6.6}{7.9}\selectfont}

% A mark at the start of a listing line, for labels set beside the listing.
% Used through the listing's escapechar; the labels are drawn by an overlay
% picture after the listing, so the build needs two passes (latexmk does that).
\newcommand{\mk}[1]{\tikz[remember picture,overlay]\coordinate (#1) at (0,0.55ex);}

\tikzset{
  rawbrace/.style={decorate,decoration={brace,amplitude=4pt},hairline,
                   line width=0.8pt},
  rawlab/.style={anchor=west,font=\footnotesize,inner sep=0pt},
  box/.style={draw=hairline,line width=0.8pt,fill=codebg,inner sep=4pt,
              align=center,minimum height=0.9cm},
  arr/.style={-{Straight Barb[length=1.6mm]},muted,line width=0.7pt},
  arra/.style={-{Straight Barb[length=1.8mm]},accent,line width=0.9pt},
  arrl/.style={-{Straight Barb[length=1.8mm]},alt,line width=0.9pt},
}

\begin{document}
\shorthandoff{"}

% ------------------------------------------------------------------ title --
% 1
\begin{frame}[plain]
  \titlepagednc{Лекция 15}{Тестирование и качество}{}{Дмитрий Курочкин}
\end{frame}

% ==================================================== 1. WHY AND WHICH ==
\section{1. Зачем и какие бывают тесты}

% 2
\begin{frame}[plain]
  \sectionopener{1}{Зачем и какие бывают тесты}{Код работает --- и код, про
    который \hi{известно}, что он работает.}
\end{frame}

% 3
\begin{frame}[plain]
  \vote{Вы поменяли одну строку в сервисе заказов. Откуда вы знаете, что
    ничего не сломалось?}
\end{frame}

% 4
\begin{frame}{Что такое тест}
  \slidetop
  \claim{Тест --- код, который \hi{вызывает ваш код} и проверяет результат.}

  \versus{Автоматический тест}{%
    \begin{cbul}{accent}
      \item выполняется каждый раз
      \item бесплатно
      \item не пропускает пунктов
    \end{cbul}}{Проверка руками}{%
    \begin{cbul}{alt}
      \item выполняется один раз
      \item минуты на каждую правку
      \item «всё проверил» --- не всё
    \end{cbul}}

  \slidegap
  \aside{Чек-лист пилота, который читает машина. Пункты пишете вы: чего в списке
    нет, то не проверено.}
\end{frame}

% 5
\begin{frame}{Зачем}
  \slidetop
  \claim{Не страх сломать, а \hi{выгода}.}

  {\renewcommand{\arraystretch}{1.6}
  \begin{tabular}{@{}L{4.4cm}L{8.4cm}@{}}
    \hi{\textbf{скорость изменений}} & через минуту известно, что сломалось;
      осторожность стоит дороже тестов \\
    \hi{\textbf{документация}}       & утверждение о поведении, которое проверяется
      каждый день \\
    \hi{\textbf{дизайн}}             & тест --- первый потребитель кода: неудобно
      ему --- неудобно всем \\
    \hi{\textbf{рефакторинг}}        & зелёные до, зелёные после --- поведение то же \\
  \end{tabular}}

  \slidegap
  \aside{Порядок «тест до кода» --- инструмент проектирования. Обязательно одно:
    тест есть \hi{до} слияния в главную ветку.}
\end{frame}

% 6
\begin{frame}{Пирамида}
  \slidetop
  \begin{center}
    \begin{tikzpicture}[font=\footnotesize,
      lv/.style={box,minimum height=0.82cm,font=\footnotesize\bfseries},
      ds/.style={anchor=west,inner sep=0pt,text=muted}]
      \node[lv,minimum width=5.6cm,draw=accent,line width=1.0pt] (u) at (0,0) {unit};
      \node[lv,minimum width=4.4cm] (i) at (0,0.98) {integration};
      \node[lv,minimum width=3.2cm] (e) at (0,1.96) {e2e};
      \node[lv,minimum width=2.0cm] (l) at (0,2.94) {нагрузочные};
      \node[ds] at (3.3,0)    {один класс · миллисекунды · \hi{тысячи}};
      \node[ds] at (3.3,0.98) {несколько частей вместе · секунды · сотни};
      \node[ds] at (3.3,1.96) {весь сервис снаружи · десятки секунд · десятки};
      \node[ds] at (3.3,2.94) {под трафиком · минуты · единицы};
      \draw[arrl] (-3.6,-0.35) -- (-3.6,3.3);
      \node[anchor=east,inner sep=0pt,text=alt,font=\scriptsize\bfseries,align=right]
        at (-3.8,1.47) {ЦЕНА};
    \end{tikzpicture}
  \end{center}

  \slidegap
  \aside{Чем выше --- тем \lo{медленнее, дороже и больше ложных падений}. Тест,
    который падает раз в неделю без причины, перестают читать.}
\end{frame}

% 7
\begin{frame}{Что ловит каждый уровень}
  \slidetop
  {\renewcommand{\arraystretch}{1.5}
  \begin{tabular}{@{}L{3.2cm}L{9.6cm}@{}}
    \hi{\textbf{unit}}        & ошибки в логике: условие, знак, статус. Ничего
      о том, как части соединены \\
    \textbf{integration}      & SQL, который не переводится; миграцию, которая не
      применяется; порядок middleware \\
    \textbf{e2e}              & развёртывание, конфигурацию, взаимодействие сервисов \\
    \textbf{нагрузочные}      & голодание пула, утечки соединений, N+1 --- всё, что
      видно только под трафиком \\
  \end{tabular}}

  \slidegap
  \aside{Вне пирамиды: контрактные тесты между сервисами, дымовые --- после
    выкладки, тесты устойчивости --- уронить базу нарочно.}
\end{frame}

% 8
\begin{frame}[plain]
  \vote{Что этот тест поймает, чего не поймает более дешёвый?}
\end{frame}

% 9
\begin{frame}{Уровень теста}
  \slidetop
  \claim{Тест пишется там, где ошибка \hi{живёт}.}

  {\renewcommand{\arraystretch}{1.6}
  \begin{tabular}{@{}L{5.2cm}L{7.6cm}@{}}
    расчёт скидки             & \hi{\textbf{unit}}: считается в памяти,
      интеграционный ничего не добавит \\
    запрос к базе через ORM   & \hi{\textbf{integration}}: с подменённым контекстом
      проверен вызов, а \lo{не SQL} \\
  \end{tabular}}

  \slidegap
  \hrulethin{0.42}
  \rulegap
  \aside{Покрытие --- \hi{карта}, а не оценка: сто процентов набираются тестами без
    утверждений. В проекте без тестов --- сначала самый рискованный путь.}
\end{frame}

% 10
\begin{frame}[plain]
  \keyline{Тест --- проверка, которая выполняется бесплатно каждый раз.\\
    Пишите его на самом дешёвом уровне, где ошибка видна.}
\end{frame}

% ============================================================ 2. TOOLS ==
\section{2. Инструменты}

% 11
\begin{frame}[plain]
  \sectionopener{2}{Инструменты}{Фреймворки отличаются \hi{моделью изоляции} и
    атрибутами.}
\end{frame}

% 12
\begin{frame}{Три фреймворка}
  \slidetop
  {\renewcommand{\arraystretch}{1.5}
  \begin{tabular}{@{}L{2.4cm}L{5.4cm}L{5.0cm}@{}}
    & \colhead{accent}{экземпляр класса} & \colhead{alt}{атрибуты} \\
    \hi{\textbf{xUnit}} & \hi{новый на каждый тест}: конструктор и
      \texttt{Dispose} & \texttt{Fact}, \texttt{Theory} \\
    \textbf{NUnit}      & один на класс                 & \texttt{SetUp},
      \texttt{TestCase} \\
    \textbf{MSTest}     & встроен в инструменты         & \texttt{TestMethod},
      \texttt{DataRow} \\
  \end{tabular}}

  \slidegap
  \hrulethin{0.42}
  \rulegap
  \aside{В xUnit поле не переживает тест: зависимость между тестами
    \hi{невозможна} случайно. Тяжёлая подготовка --- фикстура, один раз на класс.
    Дальше всё на xUnit.}
\end{frame}

% 13
\begin{frame}[fragile]{Тестовый проект}
  \slidetop
  \claim{Обычный проект, обычные пакеты, команда \hi{\texttt{test}}.}

\begin{lstlisting}[style=bash,basicstyle=\codesm,escapechar=|]
|\mk{t1}|dotnet new xunit --name HelloWorld.Tests
|\mk{t2}|dotnet add HelloWorld.Tests reference HelloWorld
|\mk{t3}|dotnet test
|\mk{t4}|dotnet test --filter "FullyQualifiedName~OrderTests"
\end{lstlisting}
  \begin{tikzpicture}[remember picture,overlay,every node/.style={rawlab}]
    \node at ([xshift=9.6cm]t1) {проект с тестами};
    \node at ([xshift=9.6cm]t2) {ссылка на тестируемый};
    \node[text=accent] at ([xshift=9.6cm]t3) {найти, выполнить, отчёт};
    \node at ([xshift=9.6cm]t4) {подмножество};
  \end{tikzpicture}

  \slidegap
  \aside{Обычно проектов два: unit --- на домен, секунды, после каждого сохранения;
    интеграционный --- на веб-проект, с контейнерами, перед слиянием.}
\end{frame}

% ================================================= 3. ANATOMY OF A TEST ==
\section{3. Как устроен тест}

% 14
\begin{frame}[plain]
  \sectionopener{3}{Как устроен тест}{Тест --- микропрограмма: вход, действие,
    \hi{одно} утверждение.}
\end{frame}

% 15
\begin{frame}[fragile]{Три абзаца}
  \slidetop
  \claim{Подготовить, выполнить, проверить --- \hi{видно} в теле.}

\begin{lstlisting}[style=csharp,basicstyle=\codexs,escapechar=|]
public sealed class OrderTests
{
    [Fact]
|\mk{a1}|    public void MarkPaid_FromNew_ChangesStatusToPaid()
    {
|\mk{a2}|        var order = new Order(new Email("a@b.c"));

|\mk{a3}|        order.MarkPaid();

|\mk{a4}|        Assert.Equal(OrderStatus.Paid, order.Status);
    }
}
\end{lstlisting}
  \begin{tikzpicture}[remember picture,overlay,every node/.style={rawlab}]
    \node[text=accent] at ([xshift=9.6cm]a1) {что --- при чём --- что ждём};
    \node at ([xshift=9.6cm]a2) {подготовка};
    \node at ([xshift=9.6cm]a3) {действие};
    \node at ([xshift=9.6cm]a4) {утверждение};
  \end{tikzpicture}

  \vspace{0.4em}
  \aside{Имя читается без тела: упал --- и уже понятно, что сломалось.}
\end{frame}

% 16
\begin{frame}[fragile]{Как читать падение}
  \slidetop
  \claim{Ожидаемое --- \hi{первым}, фактическое --- вторым.}

\begin{lstlisting}[style=csharp,basicstyle=\codesm]
Assert.Equal(order.Status, OrderStatus.New);   // actual first, expected second
\end{lstlisting}

\begin{lstlisting}[style=term,basicstyle=\termsm]
$ dotnet test
...
  Failed HelloWorld.Tests.OrderTests.MarkPaid_FromNew_ChangesStatusToPaid [7 ms]
  Error Message:
   Assert.Equal() Failure: Values differ
Expected: Paid
Actual:   New
...
\end{lstlisting}

  \aside{Здесь аргументы перепутаны, и сообщение \lo{говорит наоборот}: заказ
    на самом деле \texttt{Paid}, а ждали \texttt{New}.}
\end{frame}

% 17
\begin{frame}[fragile]{Одна логика, три входа}
\begin{lstlisting}[style=csharp,basicstyle=\codexs,escapechar=|]
|\mk{b1}|[Theory]
|\mk{b2}|[InlineData(OrderStatus.Paid)]
[InlineData(OrderStatus.Shipped)]
|\mk{b3}|[InlineData(OrderStatus.Cancelled)]
public void MarkPaid_FromNonNew_Throws(OrderStatus initial)
{
|\mk{b4}|    var order = OrderFactory.WithStatus(initial);

|\mk{b5}|    Assert.Throws<InvalidTransitionException>(() => order.MarkPaid());
}
\end{lstlisting}
  \begin{tikzpicture}[remember picture,overlay,every node/.style={rawlab}]
    \node[text=accent] at ([xshift=10.6cm]b1) {тест с параметром};
    \draw[rawbrace] ([xshift=10.35cm,yshift=0.35em]b2) -- ([xshift=10.35cm,yshift=-0.3em]b3);
    \node at ([xshift=10.6cm]$(b2)!0.5!(b3)$) {три строки в отчёте};
    \node at ([xshift=10.6cm]b4) {фабрика};
    \node at ([xshift=10.6cm]b5) {исключение};
  \end{tikzpicture}

\begin{lstlisting}[style=term,basicstyle=\termsm]
$ dotnet test --logger "console;verbosity=normal"
...
  Passed HelloWorld.Tests.OrderTests.MarkPaid_FromNonNew_Throws(initial: Cancelled) [5 ms]
  Passed HelloWorld.Tests.OrderTests.MarkPaid_FromNonNew_Throws(initial: Shipped) [< 1 ms]
  Passed HelloWorld.Tests.OrderTests.MarkPaid_FromNonNew_Throws(initial: Paid) [< 1 ms]
  Passed HelloWorld.Tests.OrderTests.MarkPaid_FromNew_ChangesStatusToPaid [5 ms]
...
\end{lstlisting}
\end{frame}

% 18
\begin{frame}[plain]
  \vote{В тесте пять утверждений. Он упал. Что сломалось?}
\end{frame}

% 19
\begin{frame}{Одна причина упасть}
  \slidetop
  \versus{Пять тестов}{%
    \begin{cbul}{accent}
      \item одно имя --- одна вещь
      \item упал ровно тест про дату оплаты
      \item остальные четыре выполнились
    \end{cbul}}{Пять утверждений}{%
    \begin{cbul}{alt}
      \item одно имя --- пять причин
      \item какое из пяти --- в отчёте
      \item после первого остальные не выполнились
    \end{cbul}}

  \slidegap
  \hrulethin{0.42}
  \rulegap
  \aside{Правило не про число вызовов \texttt{Assert}: «не пустой и содержит
    элемент» --- \hi{одна} причина, два утверждения.}
\end{frame}

% 20
\begin{frame}[plain]
  \vote{Надо ли тестировать приватные методы?}
\end{frame}

% 21
\begin{frame}{Признаки плохого теста}
  \slidetop
  {\renewcommand{\arraystretch}{1.5}
  \begin{tabular}{@{}L{4.6cm}L{8.2cm}@{}}
    \lo{\textbf{приватное через рефлексию}} & проверяйте то, что видно снаружи;
      сложный приватный метод --- спрятанный класс \\
    \lo{\textbf{подготовка на тридцать строк}} & длинный конструктор или нужна
      фабрика тестовых данных \\
    \lo{\textbf{логика внутри}}      & \texttt{if}, цикл, расчёт ожидаемого ---
      тест и код неправы согласованно \\
    \lo{\textbf{зависит от машины}}  & культура, часовой пояс, файл на диске \\
  \end{tabular}}

  \slidegap
  \aside{Ожидаемое значение --- \hi{константа}, посчитанная вручную. Тест читается
    как предложение, а не как программа.}
\end{frame}

% 22
\begin{frame}[fragile]{Подстановки}
  \vspace{0.3em}
\begin{lstlisting}[style=csharp,basicstyle=\codexxs,escapechar=|]
[Fact]
public async Task Confirm_PaidOrder_SendsEmailToCustomer()
{
|\mk{m1}|    var repository = Substitute.For<IOrderRepository>();
    var notifier = Substitute.For<INotifier>();
|\mk{m2}|    var time = new FakeTimeProvider(
        new DateTimeOffset(2026, 9, 13, 10, 0, 0, TimeSpan.Zero));
|\mk{m3}|    repository.FindAsync(42, Arg.Any<CancellationToken>())
        .Returns(OrderFactory.Paid(id: 42, email: "a@b.c"));
|\mk{m4}|    var service = new OrderService(repository, notifier, time);

    await service.ConfirmAsync(42, CancellationToken.None);

|\mk{m5}|    await notifier.Received(1).SendAsync(
        "a@b.c", Arg.Any<string>(), Arg.Any<CancellationToken>());
}
\end{lstlisting}
  \begin{tikzpicture}[remember picture,overlay,every node/.style={rawlab}]
    \node at ([xshift=10.2cm]m1) {реализация на лету};
    \node at ([xshift=10.2cm]m2) {фиксированный момент};
    \node at ([xshift=10.2cm]m3) {заготовленный ответ};
    \node at ([xshift=10.2cm]m4) {собран руками};
    \node[text=accent] at ([xshift=10.2cm]m5) {вызван с нужным адресом};
  \end{tikzpicture}

  \slidegap
  \aside{Ни базы, ни почты, ни часов. Важно одно: письмо ушло \hi{покупателю}.}
\end{frame}

% 23
\begin{frame}{Что заменяем}
  \slidetop
  \claim{Всё недетерминированное приходит \hi{снаружи}.}

  {\renewcommand{\arraystretch}{1.45}
  \begin{tabular}{@{}L{3.6cm}L{9.2cm}@{}}
    \hi{\textbf{внешний сервис}} & медленный и ненадёжный \\
    \hi{\textbf{база}}           & это интеграционный уровень \\
    \hi{\textbf{время}}          & «сейчас» --- только через \texttt{TimeProvider}:
      в тесте его ставят и двигают \\
    \hi{\textbf{случайность}}    & генератор с фиксированным сидом \\
  \end{tabular}}

  \slidegap
  \aside{Каждая подстановка --- утверждение о зависимости, которое \lo{никто не
    проверяет}. Подстановки --- для границ, не для всего подряд.}
\end{frame}

% 24
\begin{frame}{Подстановка или реализация в памяти}
  \slidetop
  \versus{Подстановка}{%
    \begin{cbul}{accent}
      \item «что-то произошло снаружи»
      \item письмо ушло, событие опубликовано
      \item проверяется факт вызова
    \end{cbul}}{Реализация в памяти}{%
    \begin{cbul}{alt}
      \item «состояние стало таким»
      \item после сохранения заказ находится
      \item проверяется результат
    \end{cbul}}

  \slidegap
  \hrulethin{0.42}
  \rulegap
  \aside{Хранилище на словаре живёт в тестовом проекте и служит десяткам тестов.
    Проверка порядка внутренних вызовов \lo{ломается} при любом рефакторинге.}
\end{frame}

% 25
\begin{frame}[plain]
  \statement{Не управлять порядком тестов. \lo{Никогда}.}
\end{frame}

% 26
\begin{frame}{Каждый тест готовит своё}
  \slidetop
  \versus{Сам, с нуля}{%
    \begin{cbul}{accent}
      \item нужен оплаченный заказ --- создал
      \item запускается один
      \item параллельно, в любом порядке
      \item падает по своей причине
    \end{cbul}}{Сценарий по шагам}{%
    \begin{cbul}{alt}
      \item второй оплачивает «тот самый»
      \item один не запустить
      \item порядок выбирает фреймворк
      \item упал первый --- упали все
    \end{cbul}}

  \slidegap
  \hrulethin{0.42}
  \rulegap
  \aside{Статические поля и общая база тест \lo{переживают}. В общей базе ---
    уникальные данные, никакого «заказа номер один».}
\end{frame}

% 27
\begin{frame}{Асинхронные тесты}
  \slidetop
  \claim{Тест возвращает \texttt{Task} и \hi{ждёт}.}

  {\renewcommand{\arraystretch}{1.5}
  \begin{tabular}{@{}L{4.4cm}L{8.4cm}@{}}
    \hi{\textbf{\texttt{async Task}}}  & фреймворк ждёт задачу; исключение из неё
      --- падение \\
    \lo{\textbf{забытый \texttt{await}}} & тест проходит \lo{всегда}: проверять
      нечего \\
    \hi{\textbf{исключение}}           & асинхронный вариант утверждения:
      \texttt{ThrowsAsync} \\
    \hi{\textbf{отмена}}               & уже отменённый токен --- метод обязан
      бросить, а не доработать \\
  \end{tabular}}

  \slidegap
  \aside{Тест, который не ждёт задачу, --- тот же забытый \texttt{await}, что
    ловит анализатор.}
\end{frame}

% 28
\begin{frame}[plain]
  \keyline{Вход, одно действие и одна причина упасть --- с именем, которое её
    называет.\\ Подстановки --- только на границах.}
\end{frame}

% ============================================ 4. TESTING THE REAL SERVICE ==
\section{4. Тесты настоящего сервиса}

% 29
\begin{frame}[plain]
  \sectionopener{4}{Тесты настоящего сервиса}{Настоящий конвейер, настоящая
    база; подменяется только \hi{внешнее}.}
\end{frame}

% 30
\begin{frame}{Приложение в памяти теста}
  \slidetop
  \claim{Тот же файл запуска, тот же конвейер --- \hi{без сети}.}

  \begin{center}
    \begin{tikzpicture}[font=\scriptsize,
      st/.style={box,minimum width=2.2cm,minimum height=1.1cm,inner sep=1pt}]
      \node[st] (s1) at (0,0)     {\textbf{тест}};
      \node[st] (s2) at (2.75,0)  {\textbf{клиент}\\от фабрики};
      \node[st,draw=accent,line width=1.0pt] (s3) at (5.5,0) {\textbf{конвейер}\\middleware};
      \node[st,draw=accent,line width=1.0pt] (s4) at (8.25,0) {\textbf{обработчик}\\и контейнер};
      \node[st] (s5) at (11.0,0)  {\textbf{база}\\в контейнере};
      \foreach \i [evaluate=\i as \j using int(\i+1)] in {1,...,4}
        \draw[arr] (s\i) -- (s\j);
      \draw[rawbrace] ([yshift=0.5em]s3.north west) -- ([yshift=0.5em]s4.north east)
        node[midway,above=0.5em,font=\scriptsize\bfseries,text=accent]
        {ХОСТ ЦЕЛИКОМ, В ПРОЦЕССЕ ТЕСТА};
    \end{tikzpicture}
  \end{center}

  \slidegap
  \aside{Сервер не слушает порт: запросы идут от клиента, который тест держит в
    руках.}
\end{frame}

% 31
\begin{frame}[plain]
  \vote{Что в таком тесте настоящее, а что нет?}
\end{frame}

% 32
\begin{frame}{Настоящее и подменённое}
  \slidetop
  \versus{Настоящее}{%
    \begin{cbul}{accent}
      \item конвейер и маршрутизация
      \item привязка и валидация
      \item обработчик ошибок, сериализация
      \item контейнер и ваши обработчики
    \end{cbul}}{Подменили сами}{%
    \begin{cbul}{alt}
      \item внешние сервисы
      \item издатель токенов
    \end{cbul}}

  \slidegap
  \hrulethin{0.42}
  \rulegap
  \aside{В памяти, а не отдельным процессом: хост поднимается \hi{за секунду},
    тест может заменить регистрацию, точка останова срабатывает из теста.}
\end{frame}

% 33
\begin{frame}{Настоящая база}
  \slidetop
  \claim{Того же типа, что в проде, --- \hi{в контейнере}, на время тестов.}

  {\renewcommand{\arraystretch}{1.5}
  \begin{tabular}{@{}L{4.6cm}L{8.2cm}@{}}
    \lo{\textbf{подстановка хранилища}} & SQL, миграции и ORM не проверены \\
    \lo{\textbf{база в памяти, SQLite}} & другой диалект: тест врёт \\
    \hi{\textbf{контейнер}}             & случайный порт, строка подключения ---
      в конфигурацию под тем же ключом \\
  \end{tabular}}

  \slidegap
  \aside{Схема --- \hi{миграциями} с нуля, а не «создать по модели»: только так
    забытая миграция находится до выкладки.}
\end{frame}

% 34
\begin{frame}[fragile]{Фабрика приложения}
\begin{lstlisting}[style=csharp,basicstyle=\ttfamily\fontsize{6.8}{7.6}\selectfont,escapechar=|]
public sealed class ApiFactory : WebApplicationFactory<Program>, IAsyncLifetime
{
|\mk{f1}|    private readonly PostgreSqlContainer _db =
        new PostgreSqlBuilder("postgres:17-alpine").Build();
|\mk{f2}|    public Task InitializeAsync() => _db.StartAsync();
    public new async Task DisposeAsync()
    {
        await _db.DisposeAsync();
        await base.DisposeAsync();
    }
    protected override void ConfigureWebHost(IWebHostBuilder builder)
    {
|\mk{f3}|        builder.UseSetting(
            "ConnectionStrings:Orders", _db.GetConnectionString());
        builder.ConfigureTestServices(services =>
        {
|\mk{f4}|            services.RemoveAll<IRatesClient>();
            services.AddSingleton(Substitute.For<IRatesClient>());
        });
    }
}
\end{lstlisting}
  \begin{tikzpicture}[remember picture,overlay,every node/.style={rawlab}]
    \node[text=accent] at ([xshift=10.4cm]f1) {настоящий PostgreSQL};
    \node at ([xshift=10.4cm]f2) {старт --- раз на класс};
    \node at ([xshift=10.4cm]f3) {своя строка подключения};
    \node at ([xshift=10.4cm]f4) {подстановка клиента};
  \end{tikzpicture}
\end{frame}

% 35
\begin{frame}{Аутентификация в тестах}
  \slidetop
  \claim{Тоже \hi{замена регистрации}.}

  \versus{Тестовая схема}{%
    \begin{cbul}{accent}
      \item имя и роли --- из заголовка
      \item тест ставит заголовок
      \item издатель не нужен
    \end{cbul}}{Настоящие токены}{%
    \begin{cbul}{alt}
      \item токены от тестового инструмента
      \item проверяется ещё и подпись
      \item ключи в тестовой конфигурации
    \end{cbul}}

  \slidegap
  \hrulethin{0.42}
  \rulegap
  \aside{Оба законны; первая проще. Без заголовка --- \hi{401}.}
\end{frame}

% 36
\begin{frame}{Контейнер и чистка}
  \slidetop
  \claim{Контейнер --- на класс. Между тестами --- \hi{чистка}.}

  {\renewcommand{\arraystretch}{1.5}
  \begin{tabular}{@{}L{4.2cm}L{8.6cm}@{}}
    \textbf{откат транзакции}       & быстро; \lo{не работает}, если код сам
      открывает транзакции \\
    \textbf{сброс таблиц}           & медленнее, зато честнее \\
    \hi{\textbf{уникальные данные}} & чистка раз на класс --- самый быстрый
      вариант \\
  \end{tabular}}

  \slidegap
  \aside{Цена --- движок контейнеров у разработчика и на сборке. Нет движка ---
    тесты пропускаются с понятным сообщением, а не падают.}
\end{frame}

% 37
\begin{frame}[fragile]{Контрактный тест}
  \vspace{0.2em}
\begin{lstlisting}[style=csharp,basicstyle=\codexxs,escapechar=|]
public sealed class CreateOrderTests(ApiFactory factory)
    : IClassFixture<ApiFactory>
{
    [Fact]
    public async Task Post_WithNegativeQuantity_Returns422WithFieldError()
    {
        var client = factory.CreateClient();
        var body = new { customerEmail = "a@b.c",
            items = new[] { new { productId = 1, quantity = -1 } } };

|\mk{c1}|        var response = await client.PostAsJsonAsync("/orders", body);

|\mk{c2}|        Assert.Equal(
            HttpStatusCode.UnprocessableEntity, response.StatusCode);
|\mk{c3}|        var error = await response.Content
            .ReadFromJsonAsync<ErrorResponse>();
|\mk{c4}|        Assert.Contains("items[0].quantity", error!.Errors.Keys);
    }
}
\end{lstlisting}
  \begin{tikzpicture}[remember picture,overlay,every node/.style={rawlab}]
    \node[text=accent] at ([xshift=10.4cm]c1) {запрос как клиент};
    \node at ([xshift=10.4cm]c2) {статус из контракта};
    \node at ([xshift=10.4cm]c3) {тело в формате ошибки};
    \node at ([xshift=10.4cm]c4) {имя поля};
  \end{tikzpicture}
\end{frame}

% 38
\begin{frame}{Почему контракт}
  \slidetop
  \claim{Запрос, статус, тело --- и \hi{ни одного} обращения к внутренностям.}

  {\renewcommand{\arraystretch}{1.6}
  \begin{tabular}{@{}L{4.2cm}L{8.6cm}@{}}
    \hi{\textbf{это обещание}}       & его нарушение ломает клиентов; внутренности
      вы вправе переписать \\
    \hi{\textbf{проходит через всё}} & фильтр, валидатор, middleware, статус,
      имена полей --- одним тестом \\
    \hi{\textbf{документация}}       & список тестов --- список того, что сервис
      обещает \\
  \end{tabular}}

  \slidegap
  \aside{Переименовали поле --- тест упал: о ломающем изменении вы узнали
    \hi{раньше} клиентов.}
\end{frame}

% 39
\begin{frame}[plain]
  \vote{Как проверить, что заказ действительно сохранился в базу?}
\end{frame}

% 40
\begin{frame}{Через тот же контракт}
  \slidetop
  \begin{center}
    \begin{tikzpicture}[font=\footnotesize,
      tb/.style={box,minimum width=2.9cm,minimum height=1.0cm}]
      \node[tb] (t1) at (0,0)    {\texttt{POST /orders}};
      \node[tb] (t2) at (3.6,0)  {\textbf{201}\\\texttt{Location}};
      \node[tb] (t3) at (7.2,0)  {\texttt{GET} по адресу};
      \node[tb,draw=accent,line width=1.0pt] (t4) at (10.8,0)
        {\textbf{200}\\то же содержимое};
      \foreach \i [evaluate=\i as \j using int(\i+1)] in {1,...,3}
        \draw[arr] (t\i) -- (t\j);
    \end{tikzpicture}
  \end{center}

  \slidegap
  \begin{cbul}{accent}
    \item клиент делает ровно это: работает --- значит, \hi{сохранение работает}
    \item в базу заглядывают, когда наружу не видно: фоновая задача пометила
      заказы
    \item фоновую задачу тест вызывает напрямую --- таймер не ждут никогда
  \end{cbul}

  \slidegap
  \aside{Мигающий тест чинится или удаляется. Перезапуск «пока не пройдёт»
    \lo{прячет} настоящую гонку.}
\end{frame}

% 41
\begin{frame}[plain]
  \keyline{Приложение поднимается в тесте целиком, база настоящая.\\
    Проверяется контракт --- обещание, а не внутренности.}
\end{frame}

% ================================================ 5. RANDOM AND FUZZING ==
\section{5. Рандом и фаззинг}

% 42
\begin{frame}[plain]
  \sectionopener{5}{Рандом и фаззинг}{Тест придумывает входы сам, но \hi{сид}
    печатается в отчёт.}
\end{frame}

% 43
\begin{frame}{Свойство вместо примера}
  \slidetop
  \versus{Свойство}{%
    \begin{cbul}{accent}
      \item верно для \hi{любого} входа
      \item сумма заказа --- сумма позиций
      \item в JSON и обратно --- равный объект
      \item сотни входов от генератора
    \end{cbul}}{Пример}{%
    \begin{cbul}{alt}
      \item один конкретный вход
      \item две позиции по сто рублей
      \item случай, который вы вообразили
      \item три входа, и они перечислены
    \end{cbul}}

  \slidegap
  \hrulethin{0.42}
  \rulegap
  \aside{Для кода с алгоритмом: разбор, преобразования, расчёты. Фаззинг границы
    сервиса: мусор в теле --- ответ четырёхсотый, \lo{не} пятисотый.}
\end{frame}

% 44
\begin{frame}[plain]
  \vote{Тест со случайными данными упал ночью в конвейере. Утром вы его
    запустили --- он прошёл. Что делать?}
\end{frame}

% 45
\begin{frame}{Сид}
  \slidetop
  \claim{Случайность без сида --- не тест, а \lo{лотерея}.}

  \begin{cbul}{accent}
    \item сид фиксируется в начале и \hi{печатается} при падении
    \item тот же сид --- тот же вход: тест падает снова, ошибка найдена
    \item хорошая библиотека сжимает вход до минимального: не «двести элементов»,
      а «один отрицательный»
  \end{cbul}

  \slidegap
  \aside{Место в пирамиде --- рядом с unit-тестами домена: чистая логика в памяти,
    сотни входов за миллисекунды.}
\end{frame}

% ============================================== 6. PERFORMANCE TESTING ==
\section{6. Перфоманс-тестирование}

% 46
\begin{frame}[plain]
  \sectionopener{6}{Перфоманс-тестирование}{Секундомер врёт; мерит инструмент,
    и только в \hi{Release}.}
\end{frame}

% 47
\begin{frame}[plain]
  \vote{Почему нельзя мерить секундомером: засечь время, вызвать метод,
    вычесть?}
\end{frame}

% 48
\begin{frame}{Шесть причин}
  \slidetop
  {\renewcommand{\arraystretch}{1.3}
  \begin{tabular}{@{}L{4.4cm}L{8.4cm}@{}}
    первый вызов          & вы измерили \lo{компиляцию} JIT \\
    метод ещё не горячий  & вы измерили \lo{черновик} \\
    сборщик мусора        & вы измерили \lo{паузу} \\
    другой процесс        & вы измерили \lo{соседа} \\
    один вызов            & одно число, а нужно \lo{распределение} \\
    сборка в Debug        & код, которого в проде \lo{не существует} \\
  \end{tabular}}

  \slidegap
  \aside{Каждая искажает результат \hi{в разы}.}
\end{frame}

% 49
\begin{frame}[fragile]{Бенчмарк}
\begin{lstlisting}[style=csharp,basicstyle=\codexxs,escapechar=|]
|\mk{n1}|[MemoryDiagnoser]
public class ParseBench
{
|\mk{n2}|    [Params(10, 1_000)]
    public int Count;

    private string[] _lines = [];

|\mk{n3}|    [GlobalSetup]
    public void Setup() => _lines = Enumerable.Range(0, Count)
        .Select(i => $"2026-09-{i % 28 + 1:00};order-{i}").ToArray();

|\mk{n4}|    [Benchmark(Baseline = true)]
    public int Substring() => _lines.Sum(l => int.Parse(l.Substring(8, 2)));

    [Benchmark]
    public int Span() => _lines.Sum(l => int.Parse(l.AsSpan(8, 2)));
}
\end{lstlisting}
  \begin{tikzpicture}[remember picture,overlay,every node/.style={rawlab}]
    \node at ([xshift=11.3cm]n1) {диагностика памяти};
    \node at ([xshift=11.3cm]n2) {размер данных};
    \node at ([xshift=11.3cm]n3) {подготовка отдельно};
    \node[text=accent] at ([xshift=11.3cm]n4) {базовый вариант};
  \end{tikzpicture}

  \aside{\texttt{dotnet run -c Release}: отдельный процесс, прогрев. В Debug
    \lo{откажется}.}
\end{frame}

% 50
\begin{frame}{Таблица}
  \slidetop
  \claim{Среднее, ошибка, отношение к базовому, \hi{аллокации}.}

  {\renewcommand{\arraystretch}{1.3}\small
  \begin{tabular}{@{}llrrrrr@{}}
    \colhead{muted}{Method} & \colhead{muted}{Count} & \colhead{muted}{Mean} &
      \colhead{muted}{Error} & \colhead{muted}{Ratio} & \colhead{muted}{Gen0} &
      \colhead{muted}{Allocated} \\
    \texttt{Substring} & 10    & 109.72 ns & 4.221 ns & 1.01 & 0.0224 & 352 B \\
    \hi{\texttt{Span}} & 10    & 60.88 ns & 1.740 ns & 0.56 & --- & --- \\
    \texttt{Substring} & 1000  & 9,850.62 ns & 195.516 ns & 1.00 & 2.0294 & 32032 B \\
    \hi{\texttt{Span}} & 1000  & 5,228.62 ns & 103.725 ns & 0.53 & --- & --- \\
  \end{tabular}}

  \slidegap
  \aside{Разница меньше ошибки --- разницы \lo{нет}. Здесь она есть: подстрока ---
    новая строка в куче на каждый элемент, окно --- нет. Мерьте горячий путь и
    сравнивайте \hi{варианты}, а не одно число.}
\end{frame}

% 51
\begin{frame}[plain]
  \statement{В третьей лекции вы угадывали.\\ Теперь вы \hi{мерите}.}
\end{frame}

% ========================================================= 7. ANALYZERS ==
\section{7. Анализаторы}

% 52
\begin{frame}[plain]
  \sectionopener{7}{Анализаторы}{Компилятор, который знает про ваши
    \hi{типичные ошибки}.}
\end{frame}

% 53
\begin{frame}[fragile]{Три строки}
  \slidetop
  \claim{Правила инженерии --- \hi{на каждой сборке}.}

\begin{lstlisting}[style=xml,basicstyle=\codesm,escapechar=|]
<PropertyGroup>
|\mk{x1}|  <AnalysisLevel>latest-all</AnalysisLevel>
|\mk{x2}|  <TreatWarningsAsErrors>true</TreatWarningsAsErrors>
|\mk{x3}|  <EnforceCodeStyleInBuild>true</EnforceCodeStyleInBuild>
</PropertyGroup>
\end{lstlisting}
  \begin{tikzpicture}[remember picture,overlay,every node/.style={rawlab}]
    \node at ([xshift=10.3cm]x1) {все встроенные правила};
    \node[text=accent] at ([xshift=10.3cm]x2) {предупреждения --- ошибки};
    \node at ([xshift=10.3cm]x3) {стиль --- тоже при сборке};
  \end{tikzpicture}

  \slidegap
  \aside{Предупреждение, которое не ломает сборку, через неделю \lo{никто не
    видит}. Строгость каждого правила --- в \texttt{.editorconfig}, одном на
    репозиторий.}
\end{frame}

% 54
\begin{frame}[fragile]{Одно слово}
  \slidetop
\begin{lstlisting}[style=csharp,basicstyle=\codesm,escapechar=|]
|\mk{r1}|var line = reader.ReadLine();
|\mk{r2}|var line = await reader.ReadLineAsync(ct);
\end{lstlisting}
  \begin{tikzpicture}[remember picture,overlay,every node/.style={rawlab}]
    \node[text=alt] at ([xshift=7.6cm]r1) {компилируется, работает, блокирует поток};
    \node[text=accent] at ([xshift=7.6cm]r2) {то, что должно быть};
  \end{tikzpicture}

  \slidegap
\begin{lstlisting}[style=term,basicstyle=\termsm,breaklines=true]
$ dotnet build | grep -m1 CA1849
/home/user/HelloWorld/Import.cs(8,20): error CA1849: 'StreamReader.ReadLine()' synchronously blocks. Await 'StreamReader.ReadLineAsync()' instead. (https://learn.microsoft.com/dotnet/fundamentals/code-analysis/quality-rules/ca1849) [/home/user/HelloWorld/HelloWorld.csproj]
\end{lstlisting}

  \rulegap
  \aside{Модель написала правдоподобно. Тесты прошли. Под нагрузкой сервис встал.
    Анализатор ловит это \hi{до коммита}.}
\end{frame}

% 55
\begin{frame}{Что ловит и чего не ловит}
  \slidetop
  \versus{Анализатор}{%
    \begin{cbul}{accent}
      \item синхронный вызов в асинхронном методе
      \item забытый \texttt{await}
      \item не освобождённый объект
      \item двойной перебор, интерполяция в логе
    \end{cbul}}{Только тесты}{%
    \begin{cbul}{alt}
      \item неправильное условие в расчёте
      \item неверный статус в ответе
      \item SQL, который не переводится
    \end{cbul}}

  \slidegap
  \hrulethin{0.42}
  \rulegap
  \aside{Включайте в \hi{первый день}. В старом проекте --- по одному правилу.
    Подавление --- точечно и с обоснованием.}
\end{frame}

% ====================================================== 8. WHERE NEXT ==
\section{8. Куда расти дальше}

% 56
\begin{frame}[plain]
  \sectionopener{8}{Куда расти дальше}{Сборка по-настоящему --- отдельный доклад.
    \hi{CI} --- условие слияния.}
\end{frame}

% 57
\begin{frame}{Конвейер}
  \slidetop
  \claim{Хоть один шаг красный --- слияние \lo{запрещено}.}

  \begin{center}
    \begin{tikzpicture}
      \foreach \i/\t in {1/{SDK\\нужной версии},2/{восстановить\\lock-файл},
                         3/{сборка и\\анализаторы},4/{unit-тесты},
                         5/{тесты в\\контейнерах}} {
        \node[box,minimum width=1.95cm,minimum height=1.1cm,inner sep=1pt,font=\scriptsize]
          (r\i) at ({(\i-1)*2.32},0) {\t};
      }
      \node[box,draw=accent,line width=1.0pt,minimum width=1.95cm,minimum height=1.1cm,
            inner sep=1pt,font=\scriptsize\bfseries] (r6) at (11.6,0) {образ\\в реестр};
      \foreach \i [evaluate=\i as \j using int(\i+1)] in {1,...,5}
        \draw[arr] (r\i) -- (r\j);
    \end{tikzpicture}
  \end{center}

  \slidegap
  \begin{cbul}{accent}
    \item каждый шаг --- команда, которую вы \hi{уже знаете}
    \item то, что делали руками в лабораторных, --- записано один раз
    \item запрещено настройкой репозитория, а не договорённостью
  \end{cbul}
\end{frame}

% 58
\begin{frame}[plain]
  \keyline{Всё, что проверяет качество, выполняется без человека\\ и запрещает
    слияние, если красное.}
\end{frame}

% 59
\begin{frame}{Пятнадцать лекций}
  \slidetop
  \claim{У вас есть \hi{сервис} --- и понимание того, что под ним.}

  {\renewcommand{\arraystretch}{1.3}
  \begin{tabular}{@{}L{2.6cm}L{10.2cm}@{}}
    \hi{\textbf{Л1--Л2}}   & от исходника до процесса: сборка, пакеты, память \\
    \hi{\textbf{Л3--Л6}}   & язык: типы, коллекции, ресурсы, асинхронность \\
    \hi{\textbf{Л7--Л9}}   & HTTP, хост и конвейер, контракт \\
    \hi{\textbf{Л10--Л12}} & зависимости и настройки, безопасность, кэш и
      наблюдаемость \\
    \hi{\textbf{Л13--Л14}} & данные: драйвер и ORM \\
    \hi{\textbf{Л15}}      & как убедиться, что всё это работает \\
  \end{tabular}}

  \slidegap
  \aside{Мастер --- не тот, кто помнит все API. Тот, кто понимает, что под ними,
    и у кого есть \hi{инструменты}, чтобы проверить.}
\end{frame}

% 60
\begin{frame}[plain]
  \shout{Продолжайте вскрывать}
\end{frame}

\end{document}
